\documentclass[11pt,a4paper]{article}
\usepackage[margin=2.5cm]{geometry}
\usepackage{amsmath,amssymb,amsthm}
\usepackage[hidelinks]{hyperref}
\usepackage{booktabs}
\usepackage{microtype}

\newtheorem{theorem}{Theorem}
\newtheorem{proposition}[theorem]{Proposition}
\newtheorem{corollary}[theorem]{Corollary}
\theoremstyle{definition}
\newtheorem{definition}[theorem]{Definition}
\newtheorem{observation}[theorem]{Observation}
\newtheorem{conjecture}[theorem]{Conjecture}
\theoremstyle{remark}
\newtheorem{remark}[theorem]{Remark}

\newcommand{\Z}{\mathbb{Z}}
\newcommand{\Q}{\mathbb{Q}}
\newcommand{\C}{\mathbb{C}}
\newcommand{\mfld}[1]{\texttt{#1}}
\newcommand{\Gal}{\mathrm{Gal}}
\newcommand{\SU}{\mathrm{SU}}
\newcommand{\Weyl}{\mathrm{Weyl}}
\newcommand{\disc}{\mathrm{disc}}

\begin{document}

\title{Galois Groups of Cusp Shapes of Hyperbolic Flavour Manifolds\\
and Weyl Groups of $\SU(2)$, $\SU(3)$, and $\SU(4)$}

\author{M.~L.~Gentry%
\thanks{Independent Researcher, Seattle, WA, USA;
\texttt{drmlgentry@protonmail.com};
ORCID: 0009-0006-4550-2663}}
\date{\today}
\maketitle

\begin{abstract}
Three hyperbolic 3-manifolds arising in the Hyperbolic Flavour
Geometry programme have cusp shapes whose minimal polynomials
over $\Q$ have Galois groups isomorphic to Weyl groups of
Standard Model gauge factors.
The lepton manifold $M_{\mathrm{PMNS}}=\mfld{m003}(-2,3)$
(the Meyerhoff manifold) has cusped parent \mfld{m003} with
cusp shape satisfying $x^2-x+1$ ($\disc=-3$, $\Gal\cong\Z/2
\cong\Weyl(\SU(2))$).
The quark manifold $M_{\mathrm{CKM}}=\mfld{m006}(-5,2)$ has
cusped parent \mfld{m006} with cusp shape satisfying
$x^3+2x+1$ ($\disc=-59$, $\Gal\cong S_3\cong\Weyl(\SU(3))$).
The extended parent \mfld{m019} has cusp shape satisfying
$x^4-x-1$ ($\disc=-283$, $\Gal\cong S_4\cong\Weyl(\SU(4))$);
its filling \mfld{m019}$(2,1)$ is a second surgery description
of $M_{\mathrm{PMNS}}$, verified by SnapPy to 15 significant
figures.

The compositum of the splitting fields of the two cusp polynomials
corresponding to the two surgery descriptions of the Meyerhoff
manifold ($\mfld{m003}$ and $\mfld{m019}$) has Galois group
$S_4\times\Z/2\cong\Weyl(\SU(4))\times\Weyl(\SU(2))$, the Weyl
group of the gauge sector of the left-handed Pati--Salam model
$\SU(4)\times\SU(2)_L$.  The cubic cusp field of $\mfld{m006}$
(discriminant $-59$) is arithmetically independent from the
other two fields.

All computations are verified by SnapPy and Sage at 300-bit
precision; residuals are below $10^{-85}$.
\end{abstract}

\section{Introduction}

The Hyperbolic Flavour Geometry (HFG) programme~\cite{GentryUnified}
associates compact hyperbolic 3-manifolds to Standard Model (SM)
flavour parameters.
The lepton mixing matrix (PMNS) is encoded by
$M_{\mathrm{PMNS}}=\mfld{m003}(-2,3)$,
the unique minimum-volume closed orientable hyperbolic
3-manifold~\cite{GMM}.
The quark mixing matrix (CKM) is encoded by
$M_{\mathrm{CKM}}=\mfld{m006}(-5,2)$.

The cusped parents of these closed manifolds carry cusp shapes
$\tau_M\in\mathbb{H}$ that satisfy minimal polynomials over $\Q$.
In this paper we compute these polynomials for three cusped parents
(\mfld{m003}, \mfld{m006}, \mfld{m019}) and identify their Galois
groups with Weyl groups of gauge Lie algebras.

A key new result is the dual surgery identity
$M_{\mathrm{PMNS}}=\mfld{m019}(2,1)$~\cite{GentryDual}:
the Meyerhoff manifold arises as a Dehn filling of both
\mfld{m003} and the arithmetically independent cusped parent
\mfld{m019}, whose cusp Galois group $S_4\cong\Weyl(\SU(4))$
extends the correspondence to a third gauge factor.

\section{Cusp Shapes and Minimal Polynomials}

All computations use SnapPy~\cite{SnapPy} at 300-bit precision.
Polynomial identifications use \texttt{algdep} at the same
precision, with residuals below $10^{-85}$.

\subsection{The SU(2) parent: \mfld{m003}}

\begin{proposition}\label{prop:m003}
The cusp shape of \mfld{m003} is $\tau = e^{i\pi/3}$,
the Eisenstein point, satisfying $\Phi_6(x)=x^2-x+1$
with $\disc=-3$ and $\Gal(\Q(\tau)/\Q)\cong\Z/2$.
\end{proposition}

\begin{proof}
$e^{i\pi/3}$ is a primitive 6th root of unity satisfying
$(e^{i\pi/3})^2-e^{i\pi/3}+1=0$.
The splitting field is $\Q(\sqrt{-3})$ with $\Gal\cong\Z/2$.
\end{proof}

\begin{remark}
$\tau=e^{i\pi/3}$ is the unique fixed point of
$\mathrm{PSL}(2,\Z)$ of order $6$, with $j(\tau)=0$.
The closed filling \mfld{m003}$(-2,3)$ is the Meyerhoff manifold.
\end{remark}

\subsection{The SU(3) parent: \mfld{m006}}

\begin{proposition}\label{prop:m006}
The cusp shape $\tau_{\mathrm{CKM}}\approx 0.2267+1.4677i$
of \mfld{m006} satisfies the irreducible depressed cubic
$x^3+2x+1=0$, with $\disc=-59$ and $\Gal\cong S_3$.
\end{proposition}

\begin{proof}
The polynomial is Eisenstein at $p=2$, hence irreducible.
The discriminant of $x^3+px+q$ is $-4p^3-27q^2=-32-27=-59$,
a negative prime, not a square in $\Q$.
For a cubic, $\Gal\cong S_3$ iff $\disc$ is not a square
in $\Q$~\cite{Lang}.
\end{proof}

\subsection{The SU(4) parent: \mfld{m019}}

\begin{proposition}\label{prop:m019}
The cusp shape of \mfld{m019} satisfies the irreducible
quartic $x^4-x-1=0$, with $\disc=-283$ and $\Gal\cong S_4$.
\end{proposition}

\begin{proof}
The polynomial $x^4-x-1$ is irreducible over $\Q$ (no rational
roots; checked by \texttt{factor} in Sage).
Its discriminant is $-283$, a negative prime.
The Galois group, computed by Sage via \texttt{GaloisGroup},
is the transitive group of degree~4 and order~24,
which is $S_4$.
\end{proof}

\begin{proposition}[Dual surgery~\cite{GentryDual}]\label{prop:dual}
$\mfld{m003}(-2,3) = \mfld{m019}(2,1) = M_{\mathrm{PMNS}}$.
\end{proposition}

\begin{proof}
SnapPy computes the volumes of both fillings as
$0.981368828\ldots$ to 15 significant figures.
Isometry is confirmed by \texttt{is\_isometric\_to} in SnapPy.
\end{proof}

\section{The Galois--Weyl Correspondence}

\begin{observation}[Galois--Weyl trichotomy]\label{obs:main}
\begin{enumerate}
\item $\Gal(\tau_{\mfld{m003}})\cong\Z/2\cong\Weyl(\SU(2))$
\item $\Gal(\tau_{\mfld{m006}})\cong S_3\cong\Weyl(\SU(3))$
\item $\Gal(\tau_{\mfld{m019}})\cong S_4\cong\Weyl(\SU(4))$
\end{enumerate}
\end{observation}

The Weyl group identifications are standard:
$\Weyl(\SU(n+1))\cong S_{n+1}$ (the $A_n$ root system).
Table~\ref{tab:main} summarises the correspondence.

\begin{table}[h]
\centering
\caption{Galois--Weyl correspondence for the three HFG cusped parents.}
\label{tab:main}
\begin{tabular}{llllll}
\toprule
Manifold & Poly & $\disc$ & $\Gal$ & $\Weyl(G)$ & $G$ \\
\midrule
\mfld{m003} & $x^2-x+1$ & $-3$   & $\Z/2$ & $\Weyl(\SU(2))$ & $\SU(2)$ \\
\mfld{m006} & $x^3+2x+1$ & $-59$ & $S_3$  & $\Weyl(\SU(3))$ & $\SU(3)$ \\
\mfld{m019} & $x^4-x-1$ & $-283$ & $S_4$  & $\Weyl(\SU(4))$ & $\SU(4)$ \\
\bottomrule
\end{tabular}
\end{table}

\section{Arithmetic Independence and the Pati--Salam Compositum}

\begin{proposition}\label{prop:arith}
The three discriminants $-3$, $-59$, $-283$ are pairwise
arithmetically independent:
\begin{enumerate}
\item $59\equiv2\pmod3$: the prime $59$ is inert in
  $\Q(\sqrt{-3})=\Q(\tau_{\mfld{m003}})$.
\item $283\equiv1\pmod3$: the prime $283$ splits in
  $\Q(\sqrt{-3})$ as $283=\pi\bar\pi$ where
  $\pi=19+6\omega\in\Z[\omega]$, $|\pi|^2=283$.
\item The three cusp fields are pairwise arithmetically
  independent; in particular, none is contained in the
  compositum of the other two.
\end{enumerate}
\end{proposition}

\begin{corollary}[Pati--Salam compositum]\label{cor:ps}
The compositum of the splitting fields of $x^2-x+1$ and
$x^4-x-1$ --- the two cusp polynomials corresponding to the
two surgery descriptions of $M_{\mathrm{PMNS}}$ --- has
Galois group
\[
\Gal(\mathrm{compositum}/\Q) \cong S_4\times\Z/2
\cong \Weyl(\SU(4))\times\Weyl(\SU(2)),
\]
the Weyl group of the gauge sector of the left-handed
Pati--Salam model $\SU(4)\times\SU(2)_L$~\cite{PS}.
The degree-$48$ compositum is confirmed by Sage.
The cubic field of \mfld{m006} is arithmetically independent
and not contained in this compositum.
\end{corollary}

\begin{remark}
The absence of an $S_5$-type cusped parent in a census scan
of 5000 manifolds~\cite{GentryUnified} is consistent with
the non-observation of $\SU(5)$ grand unification.
The sequence $\Z/2\subset S_3\subset S_4$ mirrors
$\SU(2)\subset\SU(3)\subset\SU(4)$.
\end{remark}

\section{Further Arithmetic Structure of \mfld{m006}}

\begin{proposition}\label{prop:coords}
The real part $\mathrm{Re}(\tau_{\mathrm{CKM}})$ satisfies
$8x^3-24x^2+28x-11=0$ with $\disc=-3776=-2^6\times59$.
The imaginary part squared satisfies
$64x^3-192x^2+144x-59=0$
with $\disc=-2^{12}\times3^6\times59$.
The prime $59=|\disc(x^3+2x+1)|$ appears in every discriminant
in the boundary arithmetic of \mfld{m006} but is absent from
the bulk invariant trace field discriminant
$\disc(\mathrm{ITF})=-11\times239\times103266719$.
Thus $59$ is a boundary prime of \mfld{m006}.
\end{proposition}

\section{Physical Interpretation as a Conjecture}

\begin{conjecture}[Arithmetic origin of gauge symmetries]\label{conj:gauge}
The SM gauge groups $\SU(2)$, $\SU(3)$, $\SU(4)$ arise
from the cusp-field arithmetic of \mfld{m003}, \mfld{m006},
\mfld{m019} respectively.
The $U(1)$ electromagnetic factor corresponds to
$\mathrm{Re}(\tau_{\mathrm{PMNS}})=-\tfrac12\in\Q$.
\end{conjecture}

Supporting evidence: exact polynomial identifications
(Propositions~\ref{prop:m003}--\ref{prop:m019}),
dual surgery (Proposition~\ref{prop:dual}),
pairwise arithmetic independence (Proposition~\ref{prop:arith}),
Pati--Salam compositum (Corollary~\ref{cor:ps}),
and boundary-prime structure (Proposition~\ref{prop:coords}).

\section{Conclusion}

The cusped parents \mfld{m003}, \mfld{m006}, \mfld{m019}
have cusp-shape Galois groups isomorphic to
$\Weyl(\SU(2))$, $\Weyl(\SU(3))$, $\Weyl(\SU(4))$.
Their cusp fields are pairwise arithmetically independent.
The compositum of the two parents of $M_{\mathrm{PMNS}}$
has Galois group $S_4\times\Z/2\cong\Weyl(\SU(4)\times\SU(2)_L)$.
The dual surgery $M_{\mathrm{PMNS}}=\mfld{m003}(-2,3)
=\mfld{m019}(2,1)$ links the two parents geometrically.

\section{The Eisenstein Cusp and the Optimal M\"obius Strip}

The cusp shape of \mfld{m003} is $\tau = e^{i\pi/3}$, the Eisenstein
point. The corresponding cusp torus has a fundamental domain that is
a rhombus with angles $60^\circ$ and $120^\circ$, giving an aspect ratio
\begin{equation}
\rho = \frac{\operatorname{Im}(\tau)}{|\operatorname{Re}(\tau)|}
= \frac{\sqrt{3}/2}{1/2} = \sqrt{3}.
\label{eq:aspect}
\end{equation}
This is the hexagonal (Eisenstein) lattice $\Z[\omega]$,
the unique densest packing of circles in the plane.

In a recent result, Schwartz~\cite{Schwartz} proved that the
minimum aspect ratio of a smoothly embedded (developable) M\"obius
strip in $\mathbb{R}^3$ is exactly $\sqrt{3}$.
That is, a paper M\"obius strip of width $w$ requires length
$\ell > \sqrt{3}\,w$; the bound is sharp and is achieved in the
limit by a strip whose cross-section approaches the Eisenstein rhombus.

\begin{observation}
The cusp torus of \mfld{m003} has aspect ratio $\sqrt{3}$,
equal to the optimal M\"obius strip bound of Schwartz.
Both are controlled by the Eisenstein lattice $\Z[\omega]$
with $\omega = e^{i\pi/3}$.
\end{observation}

The connection is geometric: the Eisenstein lattice is the unique
flat torus realising the hexagonal close-packing, and both the
M\"obius strip optimum and the cusp shape of \mfld{m003} are
determined by this extremal property.
The cusp of \mfld{m003} is thus, in a precise sense, the most
efficiently packed flat torus that can appear at the end of a
hyperbolic 3-manifold with $\Z/2$ cusp Galois group.

Furthermore, all six HFG manifolds (\mfld{m003}, \mfld{m006},
\mfld{m019}, \mfld{m178}, $M_{\mathrm{PMNS}}$, $M_{\mathrm{CKM}}$)
are amphicheiral --- each is isometric to its mirror image.
Amphicheirality is the 3-manifold analogue of the non-orientability
of the M\"obius strip: in both cases, orientation reversal is a
geometric symmetry of the object.
For hyperbolic 3-manifolds, amphicheirality forces the
Chern--Simons invariant to satisfy $\mathrm{CS}(M) \equiv 0$ or
$\tfrac{1}{2} \pmod{1}$~\cite{Kirk-Klassen}, which constrains
the strong CP phase $\theta_{\mathrm{QCD}} = 2\pi\,\mathrm{CS}(M)$
to zero or $\pi$.
The value $\mathrm{CS}(\mfld{m003}) = \tfrac{1}{4}$ is consistent
with amphicheirality mod~$\tfrac{1}{2}$, and gives
$\theta_{\mathrm{QCD}} = 0$ for the physical filling
$M_{\mathrm{PMNS}}$~\cite{GentryUnified}.

\section*{Data availability}
Scripts at \url{https://github.com/drmlgentry/hyperbolic-flavor-scan}.

\section*{Acknowledgements}
The author thanks D.~Gabai for correspondence.
During preparation of this manuscript the author used
Claude and DeepSeek as research assistants.
All mathematical results were independently verified using
SnapPy and SageMath.
The author takes full responsibility for all content.

\begin{thebibliography}{99}

\bibitem{GentryUnified}
M.~L.~Gentry,
``Hyperbolic Flavour Geometry,''
SSRN \href{https://doi.org/10.2139/ssrn.6775158}{6775158} (2026).

\bibitem{GentryDual}
M.~L.~Gentry,
``The Meyerhoff Manifold as a Dehn Filling of Two
Arithmetically Independent Cusped Parents,''
SSRN \href{https://doi.org/10.2139/ssrn.6845778}{6845778} (2026);
Proc.\ AMS submission 135580.

\bibitem{GentryDelta}
M.~L.~Gentry,
``A Peripheral Determinant Invariant for Hyperbolic
3-Manifolds with Quartic Cusp Field of Discriminant $-283$,''
SSRN \href{https://doi.org/10.2139/ssrn.6851440}{6851440} (2026);
AGT submission 260530-Gentry.

\bibitem{GMM}
D.~Gabai, R.~Meyerhoff, P.~Milley,
J.\ Amer.\ Math.\ Soc.\ \textbf{22}, 1157 (2009).

\bibitem{PS}
J.~C.~Pati and A.~Salam,
Phys.\ Rev.\ D \textbf{10}, 275 (1974).

\bibitem{Lang}
S.~Lang, \textit{Algebra}, 3rd ed.,
Springer, New York (2002), Ch.~VI.

\bibitem{Schwartz}
R.~E.~Schwartz,
``The optimal paper M\"obius band,''
arXiv:2308.12641 (2023).

\bibitem{Kirk-Klassen}
P.~Kirk and E.~Klassen,
``Chern--Simons invariants of 3-manifolds and representation
spaces of knot groups,''
Math.\ Ann.\ \textbf{287}, 343--367 (1990).

\bibitem{SnapPy}
M.~Culler, N.~M.~Dunfield, M.~Goerner, J.~R.~Weeks,
SnapPy, \url{http://snappy.computop.org}.

\end{thebibliography}

\end{document}
